\documentclass[10pt]{extarticle}
\usepackage[T1]{fontenc}
\usepackage[utf8]{inputenc}
\usepackage[french]{babel}
\frenchsetup{AutoSpacePunctuation=false}
\usepackage{newpxtext,newpxmath}
\usepackage[a4paper,top=1.0cm,bottom=0.9cm,left=1.15cm,right=1.15cm]{geometry}
\usepackage{microtype}
\usepackage{xcolor}
\usepackage{titlesec}
\usepackage[most]{tcolorbox}
\usepackage{enumitem}
\usepackage{multicol}

\definecolor{navy}{HTML}{1F3A5F}
\definecolor{grey}{HTML}{54606E}
\definecolor{green}{HTML}{2C6E49}
\definecolor{red}{HTML}{A6322E}
\definecolor{boxgray}{HTML}{F3F5F8}
\definecolor{boxgreen}{HTML}{F1F7F3}
\definecolor{boxedge}{HTML}{D9E0EA}

\setlength{\parindent}{0pt}
\setlength{\parskip}{2.6pt}
\setlength{\columnsep}{14pt}
\setlength{\emergencystretch}{2.5em}
\linespread{1.0}
\pagestyle{empty}

\titleformat{\section}{\color{navy}\bfseries}{}{0pt}{}
\titlespacing{\section}{0pt}{6pt}{1.5pt}

\newtcolorbox{these}{colback=boxgray,colframe=boxedge,boxrule=0.5pt,arc=2pt,
  left=7pt,right=7pt,top=3pt,bottom=3pt}
\newtcolorbox{acquis}{colback=boxgreen,colframe=green!45,boxrule=0.6pt,arc=2pt,
  left=6pt,right=6pt,top=2.5pt,bottom=2.5pt}
\newcommand{\code}[1]{\texttt{\footnotesize #1}}
\newcommand{\mm}[1]{\textbf{\textcolor{navy}{#1}}}

\begin{document}

\begin{center}
{\LARGE\bfseries\color{navy} Scella}\\[3pt]
{\large\color{grey}\itshape Prouver ce qu'une IA a fait, à quelqu'un qui n'a aucune raison de vous croire}\\[3pt]
{\footnotesize\color{grey} 7 octobre 2026 \; · \; un projet, une page \; · \; chaque chiffre est mesuré dans le code (commit \code{fb710b5}) ; la version technique donne les fichiers et les lignes}
\end{center}
\vspace{1pt}

\begin{these}
\mm{En une phrase.} Quand une décision passe par un modèle d'IA, Scella laisse derrière elle un dossier qu'un
tiers re-vérifie hors ligne --- sans installer Scella et sans croire l'éditeur --- et qui écrit lui-même ce qu'il ne
prouve pas.

\mm{Le problème.} Une organisation qui fait passer ses décisions par des modèles d'IA ne peut, aujourd'hui, presque
rien établir après coup devant un tiers. Ses traces se réécrivent sans laisser d'indice. Une suppression ne se voit
pas : une chaîne scellée dont on retire la fin reste parfaitement valide. Et les vérifier exige l'outil de celui
qu'on audite. L'AI Act impose pourtant aux systèmes à haut risque une journalisation automatique (art.~12, applicable
au 2~décembre~2027 aux systèmes de l'annexe~III, depuis le règlement~(UE)~2026/1744), et chaque cadre --- AI Act, RGPD, devoir de vigilance ---
réclame son propre dossier.
\end{these}

\begin{multicols}{2}

\section{Ce que fait Scella : trois gestes, un seul projet}

\mm{1 · Sceller au passage.} Chaque appel vers un modèle traverse un point de passage qui l'inscrit dans une chaîne
scellée : un sceau HMAC sur un encodage canonique, sous une clé propre à chaque client, qui ne voyage jamais avec le
dossier. Les données personnelles sont caviardées \emph{avant} le sceau, la dépense est bornée \emph{avant} le
sceau. Et un panel annoncé de trois modèles ne peut pas être prétendu avec un seul appel : chaque appel déclaré doit
trouver son propre enregistrement.

\mm{2 · Re-prouver sans croire l'éditeur.} Chaque dossier exporté embarque la \emph{source} de son vérifieur --- aucune
dépendance, 1\,159 lignes qu'un auditeur lit en entier. Le détenteur de la clé le recompile et re-prouve hors ligne.
Une ancre publiée ailleurs, qui annonce la racine et le nombre d'éléments, démasque une coupure par un simple
comptage, \emph{sans aucune clé}. En option : la signature de la racine et un horodatage tiers RFC~3161.

\mm{3 · Re-dériver le dossier réglementaire.} Le dossier n'est pas un document qu'on croit : c'est une fonction qu'on
ré-exécute, \emph{dossier = dériver(enregistrements scellés, lecture juridique datée)}. Le moteur re-prouve avant de
dériver et refuse le non-prouvé. Une seule capture se relit sous plusieurs cadres --- quatre sont câblés : AI Act, RGPD, CSDDD, LkSG. La lecture
est versionnée, hachée, et voyage avec le dossier : quand la règle change, l'ancien dossier reste interprétable.

\section{Ce qui le distingue}

La base --- chaîne scellée, vérifieur hors ligne, ancre --- se répand en 2026. Scella la tient avec une discipline
qui, elle, reste rare :
\begin{itemize}[leftmargin=1.1em,itemsep=1pt,topsep=1pt]
\item \mm{un cadre qui ne déclare aucun trou est refusé} dès sa construction : aucun règlement n'est entièrement
couvert par un journal, et le prétendre serait un faux vert ;
\item \mm{aucun champ «~conforme~» n'existe} : le dossier fournit l'évidence, il ne prononce pas la conformité ;
\item \mm{aucune revendication plus forte que sa preuve} : chacune est calculée depuis la source scellée, et tombe
visiblement, son motif nommé ;
\item \mm{chaque livrable écrit ce qu'il ne prouve pas}, et le remède à exiger.
\end{itemize}

\section{Ce qui tourne, mesuré}

Dix-huit composants Rust (92\,006 lignes, 1\,639 tests), plus le vérifieur. Des portes de contrôle qui font
volontairement tomber chaque garde : un contrôle incapable d'échouer ne prouverait rien. Sur une machine vierge,
\code{docker compose up} exécute un run réel, l'exporte et l'ancre, puis démasque une copie tronquée --- par
comptage, sans clé. La re-preuve complète se fait chez le détenteur de la clé. Sous le profil «~jonction~», le même run se relit
en dossier réglementaire. La configuration de production démarre contre une vraie base, survit à un
redémarrage et re-prouve son reçu.

\section{Ce qu'il ne prouve pas}

L'intégrité d'une suite d'événements, jamais la justesse d'une décision. Le dossier est inviolable pour quiconque
détient la clé, sans être «~sans confiance~» pour tous : l'ancre, sans clé, et l'horodatage tiers, en option, le
rendent opposable à son propre détenteur. Le refus de démarrer protège d'une configuration incomplète, pas d'un
opérateur malveillant.

\section{Pour qui}

Les organisations régulées qui font passer des décisions par des modèles d'IA --- banque, assurance,
santé, secteur public --- et ceux qui doivent ensuite les auditer.

\section{Où en est-on}

Aucun client encore. Les allers-retours vers les modèles distants sont éprouvés contre un serveur local qui imite
leur API : le chemin, pas encore le partenaire. La co-signature des nœuds de calcul est vérifiée, mais pas encore
produite.

\begin{acquis}
\mm{Ce que nous cherchons.} Un premier partenaire régulé --- banque, assurance, santé, secteur public --- pour un
pilote en mode observation : Scella enregistre et scelle, et un échec d'enregistrement ne bloque jamais une
requête.
\end{acquis}

\end{multicols}

\end{document}
